% !PRINTED 87
\weberpage{87}

\begin{center}
{\large \S.\ 25.}

\vspace{0.5em}
{\large\gesperrt{Elimination from Linear Equations.}}
\end{center}

It sometimes happens that, for a given system of linear equations, the question concerns not so much the actual determination of the unknowns as the decision whether their solution is possible, that is, the setting up of the condition equations which must hold among the coefficients if solutions, or solutions of a specified kind, are to exist at all. The setting up of these condition equations is called \gesperrt{elimination}. The solution of this problem is already contained implicitly in what precedes; but we shall return explicitly to some questions belonging here.

As in the preceding paragraph, we consider a system of \(m\) linear equations with \(n\) homogeneously occurring unknowns, and ask: when does this system have a solution in which not all the unknowns vanish? We have already seen that this is always the case if \(n>m\). But if \(n\le m\), the \gesperrt{necessary and sufficient} condition for such a solution is that all the \(n\)-rowed determinants of the matrix vanish. For if one of them does not vanish, then by \S.\,24, II the values of the unknowns are necessarily zero; whereas, if they all vanish, a number \(\nu<n\) can be found such that all \((\nu+1)\)-rowed determinants of the matrix are zero while at least one of the \(\nu\)-rowed determinants does not vanish, so that by \S.\,24, III a solution of the required kind exists.

From the matrix one can form, when \(n\le m\),
\[
        {m(m-1)\cdots(m-n+1)\over 1\cdot2\cdots n}
\]
\(n\)-rowed determinants, and this, therefore, would be the number of conditions. If \(n=m\), this number is \(=1\), and we obtain the case \S.\,24, II and, as was to be expected, one condition. In general, however, although all these conditions must be fulfilled, their number is greater than necessary, because some of them are necessary consequences of the others.
